\section{Related Work}
\label{sec:related_work}

\subsection{Model Merging and Task Arithmetic}
Model merging has emerged as a key strategy to unify multiple specialized deep neural networks into a single multi-task model without incurring the cost of joint multi-task training. The earliest formulations focused on weight-space interpolation of models sharing the same pre-trained initialization, often referred to as linear mode connectivity \cite{mirzadeh2020linear, frankle2020linear}. Landmark works like Model Soups \cite{wortsman2022model} and Robust Fine-Tuning \cite{wortsman2022robust} demonstrated that simple linear weight averaging of independent fine-tuned checkpoints can improve both generalizability and out-of-distribution robustness.

To adapt pre-trained models to diverse tasks, Ilharco et al.\ \cite{ilharco2022editing} formalized the concept of task vectors, where fine-tuned models are parameterized as additive updates $\tau_k = \theta_k - \theta_{\text{pre}}$. Linear combinations of these updates, known as Task Arithmetic, allow for multi-task model construction and task negation:
\begin{equation}
    \theta_{\text{merged}} = \theta_{\text{pre}} + \sum_{k=1}^K \lambda_k \tau_k
\end{equation}
where $\lambda_k$ is a scalar coefficient representing the importance or strength of task $k$. Despite its simplicity, Task Arithmetic is prone to destructive weight-space interference, especially when merging a larger number of tasks or tasks with opposing gradients. To mitigate this interference, several methods have been proposed to sparsify, mask, or scale task vectors before merging, such as TIES-Merging \cite{yadav2023ties}, DANGLES \cite{jin2023dangles}, and Fisher Merging \cite{matena2022merging}. However, these approaches still largely rely on uniform, heuristic, or global scaling parameters, which may not scale to highly diverse multi-task scenarios.

\subsection{Layer-wise and Adaptive Merging Paradigms}
To address the limitations of global task scaling, recent papers have moved towards finer-grained merging schemes that allow coefficients to vary across the layers of the model. The underlying architectural assumption is that different layers of deep neural networks capture distinct features---with early layers learning general visual or textual representations and deeper layers representing highly task-specific features \cite{dosovitskiy2020vit, nguyen2020wide}. Under this assumption, learning layer-wise coefficients $\lambda^l_k$ for each task $k$ and layer $l$ is claimed to be essential for resolving localized representational conflicts.

AdaMerging \cite{yang2024adamerging} pioneered this layer-wise adaptive merging strategy by setting up an unsupervised optimization objective. It optimizes the coefficient set $\Lambda = \{\lambda^l_k\}$ directly on an unlabeled target dataset by minimizing the entropy of the merged model's predictions:
\begin{equation}
    \min_{\Lambda} \mathcal{H}(P(y | X; \theta_{\text{merged}}(\Lambda)))
\end{equation}
Building upon this, SyMerge \cite{jung2025symerge} extends the test-time adaptation framework by jointly optimizing layer-wise merging coefficients alongside task-specific classification heads. SyMerge employs a self-labeling mechanism using specialized teacher models on the test-time data to guide the coefficient search. 

While these methods report impressive performance gains, they often do so using weak, uncalibrated baselines (e.g., standard Task Arithmetic with an unoptimized global scale $\lambda=0.3$) and fail to perform diagnostic tests on the learned parameters. As critical methodologists, we argue that the reported SOTA improvements may be a result of basic learning-rate or scaling calibration rather than fine-grained, layer-specific coordination.

\subsection{Test-Time Adaptation and the Danger of Overfitting}
Our analysis is closely connected to the field of Test-Time Adaptation (TTA), where pre-trained models are adapted to distribution shifts or downstream tasks using unlabeled test streams. Popular TTA methods like Tent \cite{wang2021tent} and other entropy-minimization strategies \cite{liang2023comprehensive, goyal2023fly} modify model parameters (e.g., batch normalization statistics or specialized adapters) at test time.

While TTA is highly flexible, it carries a significant danger of transductive overfitting \cite{ash2020deep}. Because the model parameters are optimized directly on the evaluation test set, the observed performance may fail to generalize to new, unseen samples from the same domain. This risk is amplified when the number of tuned parameters is large relative to the size of the test set. By optimizing layer-wise coefficients (e.g., $L \times K$ parameters) on a small set of target samples, layer-wise model merging is highly susceptible to transductive overfitting. This further highlights the need for simple, low-parameter baselines like a single task-wise scalar (only $K$ parameters total).

\subsection{Representational Similarity Analysis}
To understand whether the physical weight changes in model merging translate to consistent representational behaviors, we utilize representational similarity metrics from the neural network interpretability literature. Classic tools like Canonical Correlation Analysis (CCA) \cite{morcos2018insights} and Singular Vector CCA (SVCCA) \cite{raghu2017svcca} have been used to analyze representations across architectures. 

In this work, we use linear Centered Kernel Alignment (CKA) \cite{kornblith2019similarity}, which has been shown to be more robust than CCA in identifying representational correspondences between different networks. Let $X \in \mathbb{R}^{N \times D}$ and $Y \in \mathbb{R}^{N \times D}$ be the activation matrices of two different layers on the same $N$ inputs, and let $K_1 = XX^\top, K_2 = YY^\top$ be their Gram matrices. The linear CKA is defined as:
\begin{equation}
    \text{CKA}(K_1, K_2) = \frac{\text{HSIC}(K_1, K_2)}{\sqrt{\text{HSIC}(K_1, K_1) \text{HSIC}(K_2, K_2)}}
\end{equation}
where $\text{HSIC}$ is the Hilbert-Schmidt Independence Criterion. Despite its widespread use in network analysis, representational similarity metrics have been rarely applied in the model merging literature to validate whether learned merging schedules preserve the actual feature hierarchies of the original task experts. We bridge this gap by conducting the first systematic CKA analysis of merged models under different parameter treatments.
